% ============================================================
%  MODÈLE DE FORMULAIRE — 2 colonnes
%  Lecture : colonne gauche de haut en bas, puis colonne droite
%  de haut en bas, puis page suivante.
%  Titres : bleu souligné bleu | Sous-chapitres : bande grise
% ============================================================
\documentclass[9pt]{extarticle}

% ---------- Langue, police, encodage ----------
\usepackage[utf8]{inputenc}
\usepackage[T1]{fontenc}
\usepackage[french]{babel}
\usepackage{lmodern}

% ---------- Mise en page ----------
\usepackage[a4paper, margin=1cm]{geometry}
\usepackage{multicol}          % les 2 colonnes
\setlength{\columnsep}{0.5cm}  % espace entre les colonnes
\setlength{\parindent}{0pt}
\pagestyle{empty}              % pas de numéro de page

% ---------- Maths, couleurs, tableaux ----------
\usepackage{amsmath, amssymb}
\usepackage{xcolor}
\usepackage{tabularx, array}
\usepackage{titlesec}

% ---------- Couleurs ----------
\definecolor{titrebleu}{RGB}{0,70,160}
\definecolor{bandegrise}{gray}{0.82}

% ---------- Surlignage bleu (effet stabilo) ----------
\definecolor{stabibleu}{RGB}{204,233,248}
\newcommand{\stabibleu}[1]{%
  \begingroup\setlength{\fboxsep}{1.5pt}%
  \colorbox{stabibleu}{#1}\endgroup}

% ---------- TITRE (chapitre) : texte normal, surligné bleu, centré ----------
\titleformat{\section}[block]
  {\centering\large\normalfont\large}
  {}{0pt}
  {\stabibleu}
\titlespacing*{\section}{0pt}{8pt}{6pt}

% ---------- SOUS-CHAPITRE : bande grise sur toute la colonne ----------
%  (une colonne = la moitié de la page)
\newcommand{\bandegrise}[1]{%
  \colorbox{bandegrise}{%
    \parbox{\dimexpr\columnwidth-2\fboxsep\relax}{\strut#1}}}
\titleformat{\subsection}[block]
  {\normalsize\bfseries\color{black}}
  {}{0pt}
  {\bandegrise}
\titlespacing*{\subsection}{0pt}{6pt}{3pt}

% ---------- Sous-sous-chapitre (optionnel) : gras simple ----------
\titleformat{\subsubsection}[block]{\small\bfseries}{}{0pt}{}
\titlespacing*{\subsubsection}{0pt}{4pt}{2pt}

% ---------- Tableau de formules prêt à l'emploi ----------
%  Usage :
%   \begin{formules}
%     Intitulé & $formule$ & unité & remarques \\ \hline
%   \end{formules}
\renewcommand{\arraystretch}{1.4}
\newcolumntype{C}[1]{>{\centering\arraybackslash}p{#1}}
\newcolumntype{Y}{>{\raggedright\arraybackslash}X}
\newenvironment{formules}
  {\par\small\noindent
   \tabularx{\columnwidth}{|C{0.18\columnwidth}|C{0.30\columnwidth}|C{0.08\columnwidth}|Y|}
   \hline
   \textbf{Intitulé} & \textbf{Formule} & \textbf{Unité} & \textbf{Variables / remarques} \\ \hline}
  {\endtabularx\par\medskip}

% ============================================================
\begin{document}
\begin{multicols*}{2}   % multicols* = remplit gauche puis droite (pas d'équilibrage)

% ------------------------------------------------------------
% Un fichier par chapitre dans le dossier chapitres/
% Pour ajouter un chapitre : créer chapitres/chapitreN.tex et ajouter la ligne \input ci-dessous
\section{Titre du chapitre 1}

\subsection{Sous-chapitre 1.1}
\begin{formules}
  Exemple & $F = ma$ & N & $m$ : masse ; $a$ : accélération \\ \hline
  Exemple 2 & $E = \dfrac{1}{2} m v^2$ & J & $v$ : vitesse \\ \hline
\end{formules}

\subsection{Sous-chapitre 1.2}
\begin{formules}
  Intitulé & $\displaystyle \oint \vec E \cdot d\vec S = \frac{Q}{\varepsilon_0}$ & -- & Remarque \\ \hline
\end{formules}

% Texte libre possible aussi :
\subsubsection{Remarque}
Du texte, des listes ou des équations peuvent être mis
directement sous un sous-chapitre.
\[ \vec F = q\,\vec E \]

\section{Titre du chapitre 2}

\subsection{Sous-chapitre 2.1}
\begin{formules}
  Intitulé & $a = b$ & u & ... \\ \hline
\end{formules}

\subsection{Sous-chapitre 2.2}
\begin{formules}
  Intitulé & $a = b$ & u & ... \\ \hline
\end{formules}

% Pour forcer le passage à la colonne suivante : \columnbreak
% Pour forcer une nouvelle page : \newpage

\end{multicols*}
\end{document}
